\section{Campaign Analysis}
Firstly, the AFL++ status screen of this campaign shown on Figure \ref{fig:status_fork}. The stability is at $100\%$, so the harness is deterministic, which is what we want. The corpus count it at $1828$ and as we started with $80$ images so the fuzzer found $1748$ interesting elements on its own. During the campaign the map reached $33.26\%$, This means that around 2/3 of the library has never been executed. $57$ cycles have been done, so it means that the fuzzer has gone through its entire corpus of inputs $57$ complete times, mutating and testing each one.\\\\
As visible on Figure \ref{fig:fork_edges}, one can see that the number of edges increases rapidly and reaches a first plateau after around $1\;\mathrm{h}$ of execution somewhere between $3200$ and $3300$ edges. Surprisingly, after around $9\;\mathrm{h}\;30\;\mathrm{min}$, a  new region has been discovered where new edges are rapidly found until a second plateau is reached at a bit more than 3400 edges. At the end of the campaign, the last new find was 9 minutes before, so even though fuzzer was slowing down, it was still finding new edges from time to time. If we ran longer, we could have expected some new edges to be found but except if another new region was still accessible but not found, most of the edges that could be discovered with the current harness, has already been found. So with the large number of cycles and the shape of the curve, the campaign has surely reached saturation. Additionally, 238 crashes were saved during the campaign, which will be analyzed in the next section.